\documentclass[11pt]{article}

\usepackage[T1]{fontenc}
\usepackage[utf8]{inputenc}
\usepackage{lmodern}
\usepackage[margin=1.1in]{geometry}
\usepackage{amsmath,amssymb,amsthm,mathtools}
\usepackage{microtype}
\usepackage{xcolor}
\usepackage{array}
\usepackage{enumitem}
\usepackage{titlesec}
\usepackage{hyperref}

\titleformat{\section}
  {\normalfont\LARGE\bfseries}{\thesection}{1em}{}
\titleformat{\subsection}
  {\normalfont\Large\bfseries}{\thesubsection}{1em}{}

\hypersetup{
  pdftitle={Height-One Fixed-Ring Descent: Localization and Formal
    Unramifiedness},
  pdfauthor={Chloe van der Vlugt},
  colorlinks=true,
  linkcolor=blue!55!black,
  citecolor=green!40!black,
  urlcolor=blue!60!black
}

\newtheorem{theorem}{Theorem}[section]
\newtheorem{proposition}[theorem]{Proposition}
\newtheorem{corollary}[theorem]{Corollary}
\theoremstyle{remark}
\newtheorem*{remark}{Remark}

\DeclareMathOperator{\Spec}{Spec}
\DeclareMathOperator{\Frac}{Frac}
\DeclareMathOperator{\Aut}{Aut}
\DeclareMathOperator{\htop}{ht}

\newcommand{\kq}[1]{\kappa(#1)}
\newcommand{\m}{\mathfrak m}
\newcommand{\p}{\mathfrak p}
\newcommand{\q}{\mathfrak q}
\newcommand{\leanpath}[1]{\texttt{\nolinkurl{#1}}}
\newcommand{\leansrc}[1]{%
  \leanpath{HodgeConjecture/Boundary/Normalization/#1}}

\newenvironment{leanfiles}
  {\par\medskip\begingroup\footnotesize\raggedright\noindent
   \textit{Lean correspondence.}\ }
  {\par\endgroup\medskip}

\title{Height-One Fixed-Ring Descent:\\
  Localization and Formal Unramifiedness}
\author{Chloe van der Vlugt}
\date{Draft of September 10, 2026}

\begin{document}

\maketitle

\begin{abstract}
Let a finite group \(G\) act by algebra automorphisms on a domain \(S\),
and let \(I\leq H\leq G\).  For a height-one prime \(Q\subset S\), write
\(\q_J=Q\cap S^J\).  We prove that invariants of finite-group actions
satisfy going down, that every \(\q_J\) again has height one, and that
\((S^J)_{\q_J}\) is a discrete valuation ring when \(S^J\) is Noetherian
and normal.  We then construct the canonical local injection

\[
  (S^H)_{\q_H}\longrightarrow (S^I)_{\q_I}
\]

and use it to descend formal unramifiedness at the selected prime from
\(S^I\) to \(S^H\).  In a finite normal-closure normalization, this gives
an exact endpoint based over the literal full invariant ring \(S^G\).
The identification of valuation inertia with prime inertia is also proved
from the height-one localization.

The fixed-prime-inertia theorem of the Stacks Project, Tag \texttt{0BST},
is not reproved in Lean: its local unramifiedness conclusion remains an
explicit hypothesis of the machine-checked endpoint.  Mathematically,
Tag \texttt{0BST} discharges that hypothesis, and we state the resulting
corollary separately.  None of these results asserts global \'etaleness,
purity, or the Hodge conjecture.
\end{abstract}

\begin{center}
  \small Manuscript draft for author review.  No license assertion is
  made in this draft.
\end{center}

\begingroup
\small
\tableofcontents
\endgroup

\section{Introduction}
\label{sec:introduction}

Let \(S\) be a domain equipped with an action of a finite group \(G\).
For \(H\leq G\), the inclusion of fixed rings is contravariant:

\[
  I\leq H
  \qquad\Longrightarrow\qquad
  S^H\subseteq S^I.
\]

If \(Q\subset S\) is a selected prime, contraction therefore gives a
tower of local rings

\[
  (S^G)_{\q_G}\longrightarrow
  (S^H)_{\q_H}\longrightarrow
  (S^I)_{\q_I},
  \qquad \q_J=Q\cap S^J.
\]

In the intended normalization application, \(I=I_Q\) is the inertia
group of \(Q\).  Stacks Project Tag \texttt{0BST} proves that the map from
the full invariant ring \(S^G\) to the \(I_Q\)-fixed ring is
\'etale at the selected prime \cite{StacksPrimeInertiaEtale}.  The Mathlib
revision used here does not yet contain that theorem.  The formalization
therefore takes precisely its local unramifiedness conclusion as input and
proves the remaining arrows internally.

The commutative-algebra argument has four parts.

\begin{enumerate}[label=\textup{(\roman*)},leftmargin=2.6em]
  \item finite-group invariants give a special going-down argument;
  \item height one survives contraction to every subgroup-fixed ring;
  \item normality and Noetherianity turn the selected localization into a
        discrete valuation ring;
  \item formal unramifiedness descends through the canonical local map
        from \(S^H\) to \(S^I\).
\end{enumerate}

For the geometric interface, there is one additional comparison.  A
normalization diagram initially supplies the inertia group of a centered
valuation.  At height one over a Noetherian affine base, the centered
valuation ring is the localization \(S_Q\).  This identifies valuation
inertia with the ring-theoretic prime inertia occurring in
Tag \texttt{0BST}.

The purpose of this note is deliberately narrow.  It records these local
algebraic implications, their hypotheses, and their Lean verification.  It
does not turn selected-prime unramifiedness into a global finite
\'etale cover, extend a cover across a compactification, construct
a Chow projector, or prove any case of the Hodge conjecture.  The underlying
commutative-algebra ingredients are classical; the contribution recorded
here is their exact synthesis and reusable machine-checked interface.

\section{Finite invariants and height}
\label{sec:height}

Fix a commutative ring \(A\), an \(A\)-algebra \(S\) which is a domain,
and a finite group \(G\) acting on \(S\) by \(A\)-algebra automorphisms.
For \(H\leq G\), put

\[
  R_H=S^H=\{x\in S:h(x)=x\text{ for every }h\in H\}.
\]

Every \(x\in S\) is integral over \(R_H\): the orbit polynomial

\[
  \prod_{h\in H}(T-h(x))
\]

is monic and has coefficients in \(R_H\).  We also use the fact that
\(H\) acts transitively on the primes of \(S\) above any fixed prime of
\(R_H\).  Both statements are collected in Stacks Project
Tag \texttt{0BRI} \cite{StacksInvariantPrimeTransitivity}.

\begin{proposition}[Going down for a finite-group fixed ring]
\label{prop:going-down}
The integral extension \(R_H\subseteq S\) has going down.  Explicitly, if
\(\p\subseteq\q\) are primes of \(R_H\) and \(Q\subset S\) lies over
\(\q\), then there is a prime \(P\subseteq Q\) lying over \(\p\).
\end{proposition}

\begin{proof}
By lying over, choose a prime \(P_0\subset S\) above \(\p\)
\cite{StacksLyingOver}.  Apply going up to the integral extension
\(R_H\subseteq S\), the inclusion \(\p\subseteq\q\), and \(P_0\).  This
gives a prime \(Q'\supseteq P_0\) above \(\q\) \cite{StacksGoingUp}.  The
transitivity statement for invariant primes gives \(h\in H\) such that
\(h(Q')=Q\).  Then

\[
  P=h(P_0)\subseteq h(Q')=Q.
\]

Since every element of \(R_H\) is fixed by \(h\), contraction is unchanged:
\(P\cap R_H=P_0\cap R_H=\p\).  Thus \(P\) is the required lift.
\end{proof}

\begin{leanfiles}
This is formalized as
\leanpath{NormalizationInertia.fixedPoints_hasGoingDown} in
\leansrc{FixedRingLocalization.lean}.  The Lean proof performs the same
three operations: lying over, going up, and conjugation of the terminal
prime by invariant-prime transitivity.
\end{leanfiles}

\begin{proposition}[Height-one contraction]
\label{prop:height-one}
Let \(Q\subset S\) be a prime of height one.  For every \(H\leq G\),

\[
  \htop(Q\cap R_H)=1.
\]
\end{proposition}

\begin{proof}
Going down lifts every strict chain ending at \(Q\cap R_H\) to a chain of
the same length ending at \(Q\).  Hence

\[
  \htop(Q\cap R_H)\leq\htop(Q)=1.
\]

It remains to exclude height zero.  The height-one prime \(Q\) is nonzero,
because \(S\) is a domain.  In an integral extension of domains, a prime
whose contraction is zero must itself be zero.  Thus \(Q\cap R_H\neq0\).
The zero ideal is the unique minimal prime of the domain \(R_H\), so the
nonzero prime \(Q\cap R_H\) has positive height.  The two inequalities give
the result.
\end{proof}

\begin{leanfiles}
Lean separates the general chain inequality, its fixed-ring specialization,
and its height-one consequence as
\leanpath{primeHeight_under_le_of_hasGoingDown},
\leanpath{fixedPoints_primeHeight_under_le}, and
\leanpath{fixedPoints_primeHeight_under_eq_one}, respectively, all in
\leansrc{FixedRingLocalization.lean}.  Lean's prime height agrees with the
usual chain-height convention; compare Tag \texttt{00KF}
\cite{StacksPrimeHeight}.
\end{leanfiles}

\section{The selected local morphism}
\label{sec:local-map}

Suppose \(I\leq H\leq G\), and let \(Q\subset S\) be prime.  Write

\[
  \q_H=Q\cap R_H,
  \qquad
  \q_I=Q\cap R_I.
\]

The inclusion \(R_H\hookrightarrow R_I\) carries
\(R_H\setminus \q_H\) into \(R_I\setminus \q_I\).  The universal property of
localization therefore gives a canonical map

\begin{equation}
\label{eq:local-map}
  \varphi_{I,H,Q}\colon
  (R_H)_{\q_H}\longrightarrow(R_I)_{\q_I}.
\end{equation}

\begin{proposition}[Canonical local injection]
\label{prop:local-injection}
The morphism \eqref{eq:local-map} is local.  It is injective because both
fixed rings embed in the domain \(S\).
\end{proposition}

\begin{proof}
The maximal ideals of the two localizations are \(\q_H(R_H)_{\q_H}\) and
\(\q_I(R_I)_{\q_I}\).  The equality

\[
  \q_I\cap R_H=Q\cap R_I\cap R_H=Q\cap R_H=\q_H
\]

shows that the second maximal ideal contracts to the first, so
\(\varphi_{I,H,Q}\) is local.  To see injectivity, represent an element of
\((R_H)_{\q_H}\) as \(a/s\).  If its image vanishes, then some
\(t\in R_I\setminus \q_I\) satisfies \(ta=0\) in \(R_I\subseteq S\).
The domain hypothesis gives \(a=0\), and hence \(a/s=0\).
\end{proof}

\begin{leanfiles}
The map itself is
\leanpath{NormalizationInertia.fixedPointsLocalRingHom}; locality and
injectivity are
\leanpath{fixedPointsLocalRingHom_isLocalHom} and
\leanpath{fixedPointsLocalRingHom_injective}.  They are defined and proved
in \leansrc{FixedRingUnramified.lean}.
\end{leanfiles}

\begin{proposition}[The local DVR bridge]
\label{prop:dvr}
Let \(R\) be a Noetherian normal domain and let \(\q\subset R\) have height
one.  Then \(R_{\q}\) is a discrete valuation ring.
\end{proposition}

\begin{proof}
Localization preserves Noetherianity and normality
\cite{StacksNoetherianLocalization,StacksNormalLocalization}.  The
dimension of \(R_{\q}\) is the height of
\(\q\), hence is one.  Moreover \(\q\neq0\), so \(R_{\q}\) is not a field.
Thus \(R_{\q}\) is a one-dimensional Noetherian normal local domain, and
the DVR criterion applies \cite{StacksDVRCriterion}.
\end{proof}

\begin{leanfiles}
This is formalized as
\leanpath{localizationAtPrime_isDiscreteValuationRing_of_primeHeight_eq_one}
in \leansrc{FixedRingLocalization.lean}.  Its proof constructs the local
dimension bound and invokes Mathlib's equivalent characterizations of a
discrete valuation ring.
\end{leanfiles}

\section{Conditional unramified descent}
\label{sec:unramified}

For an \(A\)-algebra \(R\) and a prime \(\q\subset R\), we say that \(R\)
is \emph{unramified over \(A\) at \(\q\)} when \(R_{\q}\) is formally
unramified over \(A\).  Equivalently,

\[
  \Omega_{R_{\q}/A}=0.
\]

When \(R\) is of finite type over \(A\), this is equivalent to the familiar
local criterion: if \(\p=\q\cap A\), then
\(\kq{\q}/\kq{\p}\) is finite separable and
\(\p R_{\q}=\q R_{\q}\)
\cite{StacksUnramifiedNecessary,StacksUnramifiedLocalCriterion}.

\begin{theorem}[Height-one fixed-ring descent]
\label{thm:fixed-ring-descent}
Let \(A\) be a Noetherian commutative ring, let \(S\) be a domain finite as an
\(A\)-module, and let a finite group \(G\) act on \(S\) by \(A\)-algebra
automorphisms.  Let \(Q\subset S\) be a height-one prime and
\(I\leq H\leq G\).  Put \(R_J=S^J\) and \(\q_J=Q\cap R_J\), and assume that
\(R_H\) is integrally closed.  If

\[
  A\longrightarrow (R_I)_{\q_I}
\]

is formally unramified, then

\[
  A\longrightarrow (R_H)_{\q_H}
\]

is formally unramified.
\end{theorem}

\begin{proof}
As an \(A\)-submodule of the finite module \(S\), the fixed ring \(R_H\)
is finite over \(A\), because \(A\) is Noetherian.  It is therefore
Noetherian
\cite{StacksNoetherianSubmodules,StacksNoetherianLocalization}.
Proposition~\ref{prop:height-one} gives
\(\htop(\q_H)=1\).  Since \(R_H\) is integrally closed,
Proposition~\ref{prop:dvr} therefore shows that

\[
  U=(R_H)_{\q_H}
\]

is a DVR\@.  Likewise \(\htop(\q_I)=1\).  Put
\(V=(R_I)_{\q_I}\).  The map \(U\to V\) is the local injection of
Proposition~\ref{prop:local-injection}.  The fixed ring \(R_I\) is a finite
\(A\)-module, so its localization \(V\) is Noetherian; moreover, \(A\to R_I\) is
finite and hence of finite type.  The local unramifiedness criterion
therefore applies.

Let \(\p=\q_H\cap A=\q_I\cap A\).  Formal unramifiedness at \(\q_I\) gives

\[
  \p V=\m_V
  \qquad\text{and}\qquad
  \kq{\q_I}/\kq{\p}\text{ is finite separable}.
\]

Locality gives \(\m_V\cap U=\m_U\); together with injectivity, this embeds
\(\kq{\q_H}\) into \(\kq{\q_I}\).  Finite separability descends to
intermediate fields, so
\(\kq{\q_H}/\kq{\p}\) is finite separable.

It remains to prove \(\p U=\m_V\cap U=\m_U\).  The ideal \(\m_V\) is nonzero
because \(\q_I\) has height one.  Hence \(\p U\neq0\), since otherwise
\(\p V=(\p U)V=0\).  Every nonzero ideal of the DVR \(U\) is a power of
\(\m_U\).  If \(\p U\subsetneq\m_U\), then
\(\p U\subseteq\m_U^2\), and extension to \(V\) gives

\[
  \m_V=\p V\subseteq\m_U^2V\subseteq\m_V^2.
\]

Thus \(\m_V=\m_V^2\).  Since \(V\) is Noetherian and local, Nakayama's
lemma would give \(\m_V=0\), a contradiction.  Hence
\(\p U=\m_U\).  The local criterion now proves that \(U\), and therefore
\(R_H\) at \(\q_H\), is formally unramified over \(A\).
\end{proof}

\begin{leanfiles}
The local descent with an explicit Dedekind hypothesis is
\leanpath{fixedRing_isUnramifiedAt_of_inertia_le_localized} in
\leansrc{FixedRingUnramified.lean}.  The theorem above, which discharges
that hypothesis from height one and normality, is
\leanpath{fixedRing_isUnramifiedAt_of_inertia_le_of_primeHeight_eq_one} in
\leansrc{FixedRingLocalization.lean}.  The structure
\leanpath{InertiaFixedRingUnramifiedAt} packages only the displayed source
premise; it does not assert that \(I\) is an inertia group.  The Lean proof
uses \leanpath{Algebra.IsUnramifiedAt.of_liesOver}, whose Dedekind-domain
statement encapsulates the local descent used above.
\end{leanfiles}

\begin{remark}[Local, not global]
Theorem~\ref{thm:fixed-ring-descent} concerns the one chosen contraction
\(\q_H\).  It does not say that \(S^H\) is globally Dedekind, that
\(S^G\to S^H\) is globally unramified, or that the associated cover is
finite \'etale everywhere.
\end{remark}

The fraction-field analogue---unramifiedness descending to a finite
intermediate extension over a DVR---is classical; see Stacks Project
Tag \texttt{0EXR} \cite{StacksUnramifiedSubextensions}.
Theorem~\ref{thm:fixed-ring-descent}
is instead phrased directly at the two selected fixed-ring localizations.

\section{Normal-closure normalization}
\label{sec:normalization}

Let \(B\) be a Noetherian domain with fraction field \(K\), let \(N/K\) be
a finite normal extension, and put

\[
  S=\operatorname{IntCl}_B(N),
  \qquad
  G=\Aut_K(N).
\]

Assume that \(S\) is finite over \(B\).  For every \(H\leq G\), the fixed
normalization ring is canonically identified with

\[
  S^H\simeq\operatorname{IntCl}_B(N^H).
\]

The right-hand side is normal in this finite-extension setting
\cite{StacksIntegralClosureNormal}.  It follows from
Propositions~\ref{prop:height-one} and~\ref{prop:dvr} that, for every
height-one \(Q\subset S\), the local ring

\[
  (S^H)_{Q\cap S^H}
\]

is a DVR\@.  This conclusion is local; the fixed ring need not have global
dimension one.

\begin{leanfiles}
The fixed-field comparison is
\leanpath{fixedFieldNormalizationEquivFixedRing} in
\leansrc{InvariantIntermediate.lean}.  Its DVR specialization is
\leanpath{fixedSubgroupNormalizationRing_localization_isDiscreteValuationRing}
in \leansrc{FixedRingLocalization.lean}.  The full-invariant base is named
\leanpath{FullInvariantNormalizationRing} in
\leansrc{PrimeInertia.lean}.  Given marked data
\leanpath{D : NormalClosureData K L N} for an additional intermediate field
\(L/K\), if \(K\) is perfect and \(B\) is integrally closed, the separate
equivalence
\leanpath{baseEquivFullInvariantNormalizationRing} identifies \(B\) with
\(S^G\); no such identification is silently assumed here.  Compare the
normal/Galois integral-closure background in Tag \texttt{0BRK}
\cite{StacksGaloisIntegralClosure}.
\end{leanfiles}

\subsection{Prime inertia and the external input}
\label{subsec:prime-inertia}

For a prime \(Q\subset S\), define

\[
  I_Q=
  \{\sigma\in G:
    \sigma(Q)=Q\text{ and }
    \sigma(x)-x\in Q\text{ for every }x\in S\}.
\]

Thus \(I_Q\) is the kernel of the action of the decomposition group on the
residue field at \(Q\).  This is exactly the subgroup used in Stacks
Project Tag \texttt{0BST} \cite{StacksPrimeInertiaEtale}.

We isolate the following premise:

\begin{equation}
\label{eq:fixed-inertia-input}
  S^G\longrightarrow
  (S^{I_Q})_{Q\cap S^{I_Q}}
  \quad\text{is formally unramified.}
\end{equation}

Tag \texttt{0BST} proves the stronger local
\'etaleness statement.  Tag \texttt{09EH} is a related DVR
specialization \cite{StacksDVRInertiaEtale}.  In the formal development,
\eqref{eq:fixed-inertia-input} is an explicit argument rather than an
axiom or a proved theorem.

\begin{theorem}[Prime-inertia endpoint]
\label{thm:prime-inertia-endpoint}
In the normalization setting above, let \(Q\subset S\) have height one and
let \(H\leq G\).  Assume \(I_Q\leq H\) and the fixed-inertia input
\eqref{eq:fixed-inertia-input}.  Then

\[
  S^G\longrightarrow (S^H)_{Q\cap S^H}
\]

is formally unramified.
\end{theorem}

\begin{proof}
Take \(A=S^G\), \(I=I_Q\), and \(R_H=S^H\) in
Theorem~\ref{thm:fixed-ring-descent}.  The ring \(S^G\) is finite over the
Noetherian ring \(B\), and hence is Noetherian.  The ring \(S\) is finite
over \(S^G\), while \(S^H\) is normal by the fixed-field normalization
comparison.  All hypotheses of that theorem are therefore satisfied.
Its source premise is exactly \eqref{eq:fixed-inertia-input}, and the
subgroup hypothesis is \(I_Q\leq H\).
\end{proof}

\begin{leanfiles}
Prime inertia is
\leanpath{normalizationPrimeInertiaGroup}, and the explicit source premise
is \leanpath{NormalizationPrimeInertiaFixedRingUnramifiedAt}, both in
\leansrc{PrimeInertia.lean}.
The exact composite is
\leanpath{normalizationFixedRing_isUnramifiedAt_of_primeInertia_le} in
\leansrc{PrimeInertiaDescent.lean}.
\end{leanfiles}

\begin{corollary}[Prime-inertia descent via Tag \texttt{0BST}]
\label{cor:prime-inertia-stacks}
In the normalization setting above, let \(Q\subset S\) have height one and
let \(H\leq G\).  If \(I_Q\leq H\), then

\[
  S^G\longrightarrow (S^H)_{Q\cap S^H}
\]

is formally unramified at the selected contraction.
\end{corollary}

\begin{proof}
Tag \texttt{0BST} says that \(S^G\to S^{I_Q}\) is \'etale at
\(Q\cap S^{I_Q}\), so the localized map in
\eqref{eq:fixed-inertia-input} is formally unramified.  Apply
Theorem~\ref{thm:prime-inertia-endpoint}.
\end{proof}

\begin{remark}[Mathematical versus formal input]
Corollary~\ref{cor:prime-inertia-stacks} invokes Tag \texttt{0BST} as a
mathematical theorem.  Because that theorem is absent from the pinned
Mathlib revision, the Lean declaration stops one arrow earlier at
Theorem~\ref{thm:prime-inertia-endpoint}, with
\eqref{eq:fixed-inertia-input} as an explicit parameter.
\end{remark}

\section{From valuation inertia to prime inertia}
\label{sec:valuation}

A normalization diagram may encode ramification by a discrete valuation
ring \(V\subset N\) centered at \(Q\), rather than directly by \(I_Q\).
Write \(I_V\) for the subgroup preserving \(V\) and acting trivially on its
residue field.  Preservation of \(V\) stabilizes its maximal ideal, so
\(Q=\m_V\cap S\) gives \(\sigma(Q)=Q\) for every \(\sigma\in I_V\).
Centeredness and triviality on the residue field also give

\[
  I_V\leq I_Q:
\]

indeed, if \(\sigma\in I_V\) and \(x\in S\), then
\(\sigma(x)-x\in\m_V\), whose contraction to \(S\) is \(Q\).

The reverse inclusion needs more than centeredness.  It follows after
identifying \(V\) with \(S_Q\).  This distinction is essential: a centered
valuation ring need not be the local ring at its center without additional
hypotheses.  The definitions of decomposition and inertia groups are
reviewed in Stacks Project Tags \texttt{0BSD} and \texttt{09EC}
\cite{StacksRamificationTheory,StacksDVRInertiaDefinition}.

\begin{proposition}[Height-one localization of the centered valuation]
\label{prop:valuation-localization}
Suppose \(S\) is the finite normalization above, \(Q\subset S\) has height
one, and \(V\) is a discrete valuation subring of \(N\) with fraction field
\(N\).  Assume that the embedding \(S\to V\) is centered at \(Q\), meaning
\(\m_V\cap S=Q\), and agrees with the given embeddings of both rings into
\(N\).
Then the centered map extends to an isomorphism

\[
  S_Q\xrightarrow{\sim}V
\]

compatible with the two embeddings into \(N\).
\end{proposition}

\begin{proof}
The finite \(B\)-algebra \(S\) is Noetherian.  As an integral closure in a
finite field extension, it is normal.  Thus \(S_Q\) is a DVR by
Proposition~\ref{prop:dvr}.  Since the maximal ideal of \(V\) contracts to
\(Q\), every element of \(S\setminus Q\) maps to a unit of \(V\), and the
centered map extends uniquely to a local injection \(S_Q\to V\).

Since \(K=\Frac(B)\) and \(N/K\) is finite, the fraction field of \(S\) is \(N\).
Thus both rings are valuation subrings of the common fraction field \(N\),
and the map is compatible with their embeddings in \(N\).  Maximality of a
valuation ring under local domination forces equality;
see the valuation-ring formulation in Tag \texttt{00I9}
\cite{StacksValuationRings}.  Hence the local map is an isomorphism.  The
factorization
of a centered map through a local ring is also the scheme-theoretic content
of Tag \texttt{01J6} \cite{StacksLocalRingFactorization}.
\end{proof}

\begin{proposition}[Equality of inertia groups]
\label{prop:inertia-equality}
Under the hypotheses of Proposition~\ref{prop:valuation-localization},

\[
  I_V=I_Q.
\]
\end{proposition}

\begin{proof}
We already have \(I_V\leq I_Q\).  Conversely, let \(\sigma\in I_Q\).
Because \(\sigma(Q)=Q\), the automorphism of \(S\) extends uniquely to an
automorphism of \(S_Q\).  For \(a/t\in S_Q\), with \(t\notin Q\), the
congruences \(\sigma(a)\equiv a\pmod Q\) and
\(\sigma(t)\equiv t\pmod Q\) show that

\[
  \frac{\sigma(a)}{\sigma(t)}
  \equiv\frac{a}{t}\pmod{QS_Q}.
\]

Thus \(\sigma\) acts trivially on the residue field of \(S_Q\).  Transport
through \(S_Q\simeq V\) shows that \(\sigma\in I_V\), proving the reverse
inclusion.
\end{proof}

\begin{leanfiles}
The compatible localization datum is
\leanpath{ValuationLocalizationAt}.  Its existence at a height-one point
is \leanpath{Diagram.hasValuationLocalizationAt_of_isNoetherianRing}; the
reverse inclusion is
\leanpath{ValuationLocalizationAt.primeInertia_le_valuationInertia}; and
the equality is
\leanpath{Diagram.valuationInertia_eq_normalizationPrimeInertiaGroup_%
of_isNoetherianRing}.
These declarations are in \leansrc{LocalizationBridge.lean}.
The direct centeredness inclusion is
\leanpath{Diagram.valuationInertia_le_normalizationPrimeInertiaGroup} in
\leansrc{PrimeInertia.lean}.  The diagram-level theorem is instantiated for
\leanpath{D : Diagram C} and
\leanpath{E : RamifiedCodimensionOnePoint}.  The structures also retain the
restriction of the valuation to \(K\) as a discrete valuation, nontrivial
valuation inertia, and the assertion that \(Q\) is a ramified height-one
point.  These hypotheses are stronger than the standalone
commutative-algebra statement.
\end{leanfiles}

\begin{corollary}[Valuation-inertia endpoint]
\label{cor:valuation-endpoint}
Under the hypotheses and notation of
Proposition~\ref{prop:valuation-localization}, let \(H\leq G\).  If
\(I_V\leq H\), then
\(S^G\longrightarrow(S^H)_{Q\cap S^H}\) is formally unramified.
\end{corollary}

\begin{proof}
Proposition~\ref{prop:inertia-equality} turns \(I_V\leq H\) into
\(I_Q\leq H\).  Apply Corollary~\ref{cor:prime-inertia-stacks}.
\end{proof}

\begin{leanfiles}
The containment conversion is
\leanpath{Diagram.normalizationPrimeInertia_le_of_valuationInertia_le_%
of_isNoetherianRing}
in \leansrc{LocalizationBridge.lean}.  The final composite is
\leanpath{NormalizationInertia.Diagram.fixedRing_isUnramifiedAt_of_%
valuationInertia_le}
in \leansrc{PrimeInertiaDescent.lean}.  That declaration is the structured
normalization-diagram specialization of the slightly more general
ring-theoretic corollary stated here.  Because Tag \texttt{0BST} is not in
the pinned Mathlib revision, the Lean declaration retains
\eqref{eq:fixed-inertia-input} as an explicit argument.
\end{leanfiles}

\section{Formal status and limitations}
\label{sec:status}

The verified dependency chain is

\[
\begin{gathered}
  \text{finite invariants}
  \Longrightarrow
  \text{going down}
  \Longrightarrow
  \htop(\q_H)=1
  \Longrightarrow
  (S^H)_{\q_H}\text{ is a DVR},\\[4pt]
  I_V\leq H
  \Longrightarrow
  I_Q\leq H,
  \qquad
  \bigl[\text{fixed-inertia unramifiedness at }\q_I\bigr]
  \Longrightarrow
  \bigl[\text{unramifiedness at }\q_H\bigr].
\end{gathered}
\]

The bracketed source statement is the explicit external input to Lean and
is supplied mathematically by Tag \texttt{0BST}.  The following claims are
not results of this note:

\begin{itemize}[leftmargin=2.2em]
  \item a Lean proof of Tag \texttt{0BST} or Tag \texttt{09EH};
  \item global unramifiedness or global \'etaleness of a fixed-ring
        cover;
  \item a purity theorem or extension across a compactification boundary;
  \item construction or extension of a Chow projector;
  \item any partial or complete proof of the Hodge conjecture.
\end{itemize}

In particular, replacing the literal source \(S^G\) by the original affine
ring \(B\) is a separate theorem.  The formalized equivalence requires marked
normal-closure data, a perfect fraction field \(K\), and an integrally closed
base \(B\); it has not been composed into the endpoint above.  Without such
hypotheses, no identification with \(B\) is asserted: \(S^G\) remains the
integral closure of \(B\) in the full fixed field \(N^G\).

The checked development uses Lean~4.24.0 and Mathlib revision
\texttt{f897ebcf72cd16f89ab4}\allowbreak
\texttt{577d0c826cd14afaafc7}.  The proof modules
cited in this note contain no \leanpath{sorry}, \leanpath{admit}, or added
axiom; the fixed-inertia statement is exposed instead as the explicit
premise \eqref{eq:fixed-inertia-input}.

The working source is at
\url{https://github.com/what-social-construct/HodgeConjecture}.  This draft
does not yet name an immutable release commit or archive; one must be
assigned before publication.

A compact challenge-facing surface appears in the repository-root files
\leanpath{Challenge.lean} and \leanpath{Solution.lean}.  The former
intentionally contains two \leanpath{sorry} placeholders for the advertised
proof holes; the checked implementations live in \leanpath{Solution.lean}
and the imported proof modules.  The surface advertises only height-one
contraction and the conditional fixed-ring descent theorem.

\section{Conclusion}
\label{sec:conclusion}

The height-one fixed-ring problem admits a clean local factorization.  A
finite group action supplies going down by combining lying over, going up,
and transitivity on primes over a fixed contraction.  This preserves height
one.  Normality and Noetherianity then make the selected fixed-ring
localization a DVR, while subgroup containment produces a canonical local
injection between the relevant localizations.  These ingredients descend
formal unramifiedness from the inertia-fixed ring to any larger subgroup
containing inertia.

For a finite normal-closure normalization over a Noetherian base, the same
height-one localization identifies centered valuation inertia with prime
inertia.  Consequently the valuation-theoretic subgroup condition used by
the normalization diagram feeds the fixed-ring descent without changing
the source ring: it remains \(S^G\).

Tag \texttt{0BST} supplies the fixed-prime-inertia premise mathematically,
so Corollaries~\ref{cor:prime-inertia-stacks}
and~\ref{cor:valuation-endpoint} follow without an additional mathematical
hypothesis.  In Lean the endpoint remains exact but conditional: all
localization and descent steps after that explicit premise are verified in
Lean~4 with Mathlib \cite{Lean4,Mathlib}.

\paragraph{Acknowledgment.}
OpenAI Codex assisted with the Lean formalization and manuscript draft.
The manuscript remains subject to author review; responsibility for its
final statements, proofs, and citations belongs to the author.

\begingroup
\small
\bibliographystyle{alpha}
\bibliography{references}
\endgroup

\end{document}
